\documentclass[12pt,a4paper]{article}

% -------------------------------------------------
% Paquetes
% -------------------------------------------------
\usepackage[utf8]{inputenc}
\usepackage[T1]{fontenc}
\usepackage[spanish,es-nodecimaldot]{babel}
\usepackage{amsmath,amssymb}
\usepackage{siunitx}
\usepackage{booktabs}
\usepackage{geometry}
\usepackage{float}
\usepackage{pgfplots}
\usepackage{caption}
\usepackage{array}

\pgfplotsset{compat=1.18}
\geometry{margin=2.5cm}

\title{\textbf{Evaluación 1}\\
Herramientas Computacionales para Ingeniería II\\
\large ``Deflexión v/s Carga''}
\author{}
\date{}

\begin{document}

\maketitle

\section{Introducción}

El presente informe analiza el comportamiento carga-deflexión de una viga
simplemente apoyada sometida a una carga puntual en el centro de su luz. El
objetivo es comparar los valores obtenidos de deflexión de forma teórico a
través de la ecuación de Euler Bernoulli con los valores sintéticos de
deflexión aportados por el docente en el enunciado del trabajo.

\section{Desarrollo}

La ecuación de Euler-Bernoulli para determinar la deflexión de una viga
simplemente apoyada en el centro de su luz es:

\begin{equation}
\delta = \frac{P L^3}{48EI}
\label{eq:euler}
\end{equation}

\noindent
\textbf{Ecuación 1:} Fórmula de Euler-Bernoulli para determinar la deflexión
en el centro de una viga de largo $L$ simplemente apoyada sometida a una
carga puntual de magnitud $P$.

Donde:
\begin{itemize}
    \item $I$: momento de inercia de la sección rectangular.
    \item $E$: módulo de la elasticidad.
\end{itemize}

Antes de emplear la fórmula teórica, es importante realizar las conversiones
pertinentes de modo que exista una coherencia dimensional. Todas las
dimensiones se convertirán a su equivalente en el sistema internacional, es
decir, para la fuerza se trabaja con N, para la presión/tensión en Pa y para
unidades de longitud en m.

\[
1\,\mathrm{GPa}=25.000.000.000\,\mathrm{Pa}
\]

Esta conversión corresponde al módulo de elasticidad.

\[
1\,\mathrm{kN}=1000\,\mathrm{N}
\]

Esta conversión corresponde a las cargas aplicadas.

Una vez se han realizado las conversiones establecidas, se calcula el
momento de inercia y próximo a ello inmediatamente se pasa a calcular la
deflexión mediante el procedimiento teórico.

\begin{equation}
I=\frac{bh^3}{12}
\label{eq:inercia}
\end{equation}

\noindent
\textbf{Ecuación 2:} Fórmula para determinar el momento de inercia respecto
a su base para una sección rectangular.

Con los valores obtenidos de la deflexión mediante la fórmula
Euler-Bernoulli, se procede a calcular la diferencia relativa respecto a los
valores sintéticos proporcionados por el docente.

\section{Resultados}

\begin{table}[H]
\centering
\caption{Resultados generales}
\label{tab:resultados}
\begin{tabular}{S[table-format=5.0]
                S[table-format=1.5]
                S[table-format=1.5]
                S[table-format=1.1]
                S[table-format=1.9]}
\toprule
{Carga [N]} & {Deflexión medida [m]} & {Deflexión teórica [m]} &
{Diferencia relativa [\%]} & {$\delta_m/P$} \\
\midrule
0     & 0.00000 & 0.00000 & 0.0 & 0.000000000 \\
5000  & 0.00026 & 0.00025 & 4.0 & 0.000000052 \\
10000 & 0.00050 & 0.00050 & 0.0 & 0.000000050 \\
15000 & 0.00076 & 0.00075 & 1.3 & 0.000000051 \\
20000 & 0.00101 & 0.00100 & 1.0 & 0.000000051 \\
25000 & 0.00128 & 0.00125 & 2.4 & 0.000000051 \\
30000 & 0.00154 & 0.00150 & 2.7 & 0.000000051 \\
35000 & 0.00177 & 0.00175 & 1.1 & 0.000000051 \\
40000 & 0.00205 & 0.00200 & 2.5 & 0.000000051 \\
\bottomrule
\end{tabular}
\end{table}

La comparación muestra que las deflexiones proporcionadas presentan una
correspondencia cercana al modelo teórico. La mayor diferencia relativa
corresponde a la carga de 5 kN, con un valor de 4\%. Para el resto de las
cargas las diferencias relativas se mantienen por debajo de este valor.

La relación carga deflexión se presenta en el siguiente gráfico:

\begin{figure}[H]
\centering
\begin{tikzpicture}
\begin{axis}[
    width=0.9\textwidth,
    height=8cm,
    xlabel={Carga [N]},
    ylabel={Deflexión [m]},
    xmin=0, xmax=45000,
    ymin=0, ymax=0.0025,
    xtick={0,5000,10000,15000,20000,25000,30000,35000,40000},
    ytick={0,0.0005,0.001,0.0015,0.002,0.0025},
    grid=major,
    legend pos=north west,
    scaled y ticks=false,
    yticklabel style={/pgf/number format/fixed},
    legend cell align=left
]
\addplot+[mark=*] coordinates {
    (0,0)
    (5000,0.00026)
    (10000,0.00050)
    (15000,0.00076)
    (20000,0.00101)
    (25000,0.00128)
    (30000,0.00154)
    (35000,0.00177)
    (40000,0.00205)
};
\addlegendentry{Deflexión medida}

\addplot+[mark=square*] coordinates {
    (0,0)
    (5000,0.00025)
    (10000,0.00050)
    (15000,0.00075)
    (20000,0.00100)
    (25000,0.00125)
    (30000,0.00150)
    (35000,0.00175)
    (40000,0.00200)
};
\addlegendentry{Deflexión teórica}
\end{axis}
\end{tikzpicture}
\caption{Gráfico Deflexión v/s Carga}
\label{fig:deflexion}
\end{figure}

Se observa un comportamiento aproximadamente lineal en ambas series, lo que
resulta consistente con la expresión de Euler-Bernoulli empleada, ya que,
manteniendo constantes $L$, $E$ e $I$ la deflexión resulta directamente
proporcional a la carga aplicada.

\section{Conclusiones}

Con verificación adicional, se puede observar que la relación entre
deflexión y carga se mantiene aproximadamente constante para los valores
analizados. Esto constituye una comprobación independiente de la tendencia
lineal esperada por el modelo.

A grandes rasgos, los resultados muestran que la formulación teórica
presenta de manera adecuada el comportamiento de los datos considerados.
Sin embargo, esto puede deberse a las condiciones específicas del problema,
ya que el modelo supone un comportamiento idealizado de la viga y los datos
proporcionados por el docente corresponden a un caso sintético proporcionado
con fines docentes.

\end{document}
